% Vietnamese translation for OI Public Library
% Source-File: IOI中国国家候选队论文/2006/黄晓愉/黄晓愉.ppt
\documentclass[11pt]{article}
\usepackage{oipl-vn}

\title{Tối ưu hóa tìm kiếm\\{\large Ghi chú theo slide}}
\author{Huang Xiaoyu}
\date{2006}

\begin{document}
\maketitle

\origtitle{Search optimization techniques}
\sourcepath{IOI China candidate team papers / 2006 / Huang Xiaoyu / Huang Xiaoyu.ppt}

\section{Phạm vi}

Bộ slide đi kèm bài viết về các kỹ thuật tối ưu tìm kiếm theo chiều sâu và
chiều rộng. Ghi chú này dựa trên bản bài viết đã dịch.

\section{Tìm kiếm theo chiều sâu}

Với bài thỏa ràng buộc, hiệu quả của DFS phụ thuộc mạnh vào thứ tự chọn biến
và thứ tự thử giá trị. Nếu điều kiện mạnh được kiểm tra sớm, cây tìm kiếm bị
cắt ngay gần gốc. Vì vậy, thay vì duyệt theo thứ tự tự nhiên của dữ liệu, nên
chọn đối tượng có ít khả năng hợp lệ nhất hoặc ảnh hưởng nhiều nhất tới phần
còn lại.

\section{Tìm kiếm theo chiều rộng}

BFS phù hợp khi cần số bước ít nhất hoặc duyệt trạng thái theo lớp. Tối ưu hóa
BFS thường xoay quanh cách mã hóa trạng thái, khử trùng lặp bằng bảng băm, và
giảm số cạnh sinh ra từ mỗi trạng thái. Nếu trạng thái có đối xứng, chuẩn hóa
trước khi đưa vào hàng đợi có thể giảm đáng kể không gian.

\section{Cắt tỉa}

Cắt tỉa cần một cận đúng. Với bài tối ưu, cận dưới hoặc cận trên giúp loại một
nhánh khi ngay cả kết quả tốt nhất trong nhánh đó cũng không thể cải thiện
đáp án hiện có. Với bài khả thi, kiểm tra ràng buộc cục bộ và suy diễn bắt
buộc giúp phát hiện mâu thuẫn sớm.

\section{Ghi nhớ}

Tối ưu tìm kiếm không chỉ là viết đệ quy nhanh. Cần thiết kế thứ tự tìm, cách
biểu diễn trạng thái, phép kiểm tra lặp, và cận cắt tỉa như một phần của cùng
một thuật toán.

\end{document}
